\documentclass[12pt]{article}
\usepackage[a3paper, margin=1in]{geometry}
\usepackage{amsmath, amssymb, amsthm, ulem}

\begin{document}

\section*{Domácí úkol 3.2}

Vypracoval: Daniel \textit{"Randál"} Ransdorf \hfill Podpis: \rule{4cm}{0.4pt}

\begin{enumerate}
  \item \textbf{Řešení}:
    Označme $a,b,c,d \in \mathbb{Z}_5$ hodnoty v matici $X$ takto:
    \[
      X = \left(\begin{array}{cc}
        a & b \\
        c & d
      \end{array}\right)
    \]
    Má platit vztah:
    \[
      \left(\begin{array}{cc}
        a & b \\
        c & d
      \end{array}\right) \cdot
      \left(\begin{array}{cc}
        2 & 1 \\
        3 & 0
      \end{array}\right) = 
      \left(\begin{array}{cc}
        2 & 1 \\
        3 & 0
      \end{array}\right) \cdot
      \left(\begin{array}{cc}
        a & b \\
        c & d
      \end{array}\right)
    \]
    Vynásobme tedy matice na obou stranách:
    \[
      \left(\begin{array}{cc}
        2a+3b & a \\
        2c+3d & c
      \end{array}\right) = 
      \left(\begin{array}{cc}
        2a+c & 2b+d \\
        3a & 3b
      \end{array}\right) 
    \]
    Z toho dostáváme čtyři rovnosti:
    \[
      \left\{
      \begin{array}{rcl}
      2a+3b&=&2a+c \\
      a&=&2b+d \\
      2c+3d&=&3a \\
      c&=&3b
      \end{array}
      \right\}
      \quad \Longrightarrow \quad
      \left\{
      \begin{array}{rcl}
      0a + 3b + 4c +0d &=& 0 \\
      1a + 3b + 0c + 4d &=& 0 \\
      2a + 0b + 2c + 3d &=& 0 \\
      0a + 2b + 1c + 0d &=& 0
      \end{array}
      \right\}
    \]
    Řešme tedy soustavu:
    \[
      \left(\begin{array}{cccc}
        0 & 3 & 4 & 0 \\
        1 & 3 & 0 & 4 \\
        2 & 0 & 2 & 3 \\
        0 & 2 & 1 & 0
      \end{array}\right) \cdot
      \left(\begin{array}{c}
        a \\ b \\ c \\ d
      \end{array}\right) = 
      \left(\begin{array}{c}
        0 \\ 0 \\ 0 \\ 0
      \end{array}\right) 
    \]
    \[
      \left(\begin{array}{cccc|c}
        0 & 3 & 4 & 0 & 0 \\
        1 & 3 & 0 & 4 & 0 \\
        2 & 0 & 2 & 3 & 0 \\
        0 & 2 & 1 & 0 & 0
      \end{array}\right)
      \overset{\text{3.ř } (\cdot 2)}{\;\sim\;}
      \left(\begin{array}{cccc|c}
        0 & 3 & 4 & 0 & 0 \\
        1 & 3 & 0 & 4 & 0 \\
        4 & 0 & 4 & 1 & 0 \\
        0 & 2 & 1 & 0 & 0
      \end{array}\right)
      \overset{\text{3.ř } (+ \text{2.ř})}{\;\sim\;}
      \left(\begin{array}{cccc|c}
        0 & 3 & 4 & 0 & 0 \\
        1 & 3 & 0 & 4 & 0 \\
        0 & 3 & 4 & 0 & 0 \\
        0 & 2 & 1 & 0 & 0
      \end{array}\right)
      \overset{\text{4.ř } (+ \text{1.ř})}{\;\sim\;}
      \left(\begin{array}{cccc|c}
        0 & 3 & 4 & 0 & 0 \\
        1 & 3 & 0 & 4 & 0 \\
        0 & 3 & 4 & 0 & 0 \\
        0 & 0 & 0 & 0 & 0
      \end{array}\right)
    \]
    \[
      \overset{\text{1.ř} \leftrightarrow \text{2.ř}}{\;\sim\;}
      \left(\begin{array}{cccc|c}
        1 & 3 & 0 & 4 & 0 \\
        0 & 3 & 4 & 0 & 0 \\
        0 & 3 & 4 & 0 & 0 \\
        0 & 0 & 0 & 0 & 0
      \end{array}\right)
      \overset{\text{2.ř } (\cdot 4)}{\;\sim\;}
      \left(\begin{array}{cccc|c}
        1 & 3 & 0 & 4 & 0 \\
        0 & 2 & 1 & 0 & 0 \\
        0 & 3 & 4 & 0 & 0 \\
        0 & 0 & 0 & 0 & 0
      \end{array}\right)
      \overset{\text{3.ř } (+ \text{2.ř})}{\;\sim\;}
      \left(\begin{array}{cccc|c}
        1 & 3 & 0 & 4 & 0 \\
        0 & 2 & 1 & 0 & 0 \\
        0 & 0 & 0 & 0 & 0 \\
        0 & 0 & 0 & 0 & 0
      \end{array}\right)
    \]
    \hspace{2cm} $c,d$ jsou volné proměnné, položme $c=t_1, d=t_2$ a vyjádřeme $a,b$:
    \[
      \;\sim\;
      \left(\begin{array}{cccc|c}
        1 & 3 & 0 & 4 & 0 \\
        0 & 2 & 1 & 0 & 0 \\
        0 & 0 & 0 & 0 & 0 \\
        0 & 0 & 0 & 0 & 0 \\ \hline
        0 & 0 & 1 & 0 & t_1 \\
        0 & 0 & 0 & 1 & t_2 \\
      \end{array}\right)
      \overset{\text{1.ř } (+ \text{6.ř})}{\;\sim\;}
      \left(\begin{array}{cccc|c}
        1 & 3 & 0 & 0 & t_2 \\
        0 & 2 & 1 & 0 & 0 \\
        0 & 0 & 0 & 0 & 0 \\
        0 & 0 & 0 & 0 & 0 \\ \hline
        0 & 0 & 1 & 0 & t_1 \\
        0 & 0 & 0 & 1 & t_2 \\
      \end{array}\right)
      \overset{\text{5.ř } (\cdot 4)}{\;\sim\;}
      \left(\begin{array}{cccc|c}
        1 & 3 & 0 & 0 & t_2 \\
        0 & 2 & 1 & 0 & 0 \\
        0 & 0 & 0 & 0 & 0 \\
        0 & 0 & 0 & 0 & 0 \\ \hline
        0 & 0 & 4 & 0 & 4t_1 \\
        0 & 0 & 0 & 1 & t_2 \\
      \end{array}\right)
      \overset{\text{2.ř } (+ \text{5.ř})}{\;\sim\;}
      \left(\begin{array}{cccc|c}
        1 & 3 & 0 & 0 & t_2 \\
        0 & 2 & 0 & 0 & 4t_1 \\
        0 & 0 & 0 & 0 & 0 \\
        0 & 0 & 0 & 0 & 0 \\ \hline
        0 & 0 & 4 & 0 & 4t_1 \\
        0 & 0 & 0 & 1 & t_2 \\
      \end{array}\right)
    \]
    \[
      \overset{\text{1.ř } (+ \text{2.ř})}{\;\sim\;}
      \left(\begin{array}{cccc|c}
        1 & 0 & 0 & 0 & 4t_1 + t_2 \\
        0 & 2 & 0 & 0 & 4t_1 \\
        0 & 0 & 0 & 0 & 0 \\
        0 & 0 & 0 & 0 & 0 \\ \hline
        0 & 0 & 4 & 0 & 4t_1 \\
        0 & 0 & 0 & 1 & t_2 \\
      \end{array}\right)
      \overset{\text{2.ř } (\cdot 3)}{\;\sim\;}
      \left(\begin{array}{cccc|c}
        1 & 0 & 0 & 0 & 4t_1 + t_2 \\
        0 & 1 & 0 & 0 & 2t_1 \\
        0 & 0 & 0 & 0 & 0 \\
        0 & 0 & 0 & 0 & 0 \\ \hline
        0 & 0 & 4 & 0 & 4t_1 \\
        0 & 0 & 0 & 1 & t_2 \\
      \end{array}\right)
      \overset{\text{5.ř } (\cdot 4)}{\;\sim\;}
      \left(\begin{array}{cccc|c}
        1 & 0 & 0 & 0 & 4t_1 + t_2 \\
        0 & 1 & 0 & 0 & 2t_1 \\
        0 & 0 & 0 & 0 & 0 \\
        0 & 0 & 0 & 0 & 0 \\ \hline
        0 & 0 & 1 & 0 & t_1 \\
        0 & 0 & 0 & 1 & t_2 \\
      \end{array}\right)
    \]
    \[
      \left(\begin{array}{c}
        a \\ b \\ c \\d
      \end{array}\right) =
      \left(\begin{array}{c}
        4t_1 + t_2 \\ 2t_1 \\ t_1 \\ t_2
      \end{array}\right) \quad , t_1,t_2 \in \mathbb{Z}_5
    \]

    Z řešení soustavy rovnic můžeme tedy dosazením $a,b,c,d$ vyjádřit všechny matice X:

    \[
      X = \left(\begin{array}{cc}
        4t_1 + t_2 & 2t_1 \\
        t_1 & t_2
      \end{array}\right) \quad , t_1,t_2 \in \mathbb{Z}_5
    \]
    
    Popřípadě můžeme matici X zapsat jako součet matic násobených parametry:

    \[
      X =  \underline{\underline{t_1\left(\begin{array}{cc}
        4 & 2 \\
        1 & 0
      \end{array}\right) +
      t_2\left(\begin{array}{cc}
        1 & 0 \\
        0 & 1
      \end{array}\right)}}
      \quad , t_1,t_2 \in \mathbb{Z}_5
    \]

\end{enumerate}

\end{document}
